\documentclass[11pt,letterpaper]{article}
\usepackage[utf8]{inputenc}
\usepackage[T1]{fontenc}
\usepackage[spanish,es-tabla,es-noshorthands]{babel}
\usepackage{lmodern}
\usepackage[margin=2.5cm]{geometry}
\usepackage{amsmath,amssymb,bm}
\usepackage{booktabs,array}
\usepackage{graphicx}
\usepackage{float}
\usepackage{enumitem}
\usepackage{tikz}
\usetikzlibrary{positioning,arrows.meta}
\usepackage{quantikz}
\usepackage[colorlinks=true,linkcolor=blue!50!black,citecolor=blue!50!black,urlcolor=blue!50!black]{hyperref}
\usepackage{caption}
\captionsetup{font=small,labelfont=bf}

\providecommand{\ket}[1]{\left|#1\right\rangle}
\providecommand{\bra}[1]{\left\langle#1\right|}
\providecommand{\braket}[2]{\left\langle#1\middle|#2\right\rangle}
\newcommand{\R}{\mathbb{R}}

\title{\textbf{Propuesta de proyecto final:}\\[2pt]
Predicción del desenlace de perros y gatos en un refugio desde su ingreso con un clasificador cuántico de \emph{data re-uploading}}
\author{%
\begin{tabular}{ll}
\toprule
Integrante & No.\ de cuenta \\ \midrule
Rojas Chimal Karla Beatriz & 320048025 \\
Calderón Gutiérrez Victor Emiliano & 423097331 \\
\bottomrule
\end{tabular}\\[10pt]
Computación Cuántica -- Proyecto Final\\
Prof.\ Naomi Itzel Reyes Granados \;\;|\;\; Ayud.\ Sebastián González Juárez}
\date{Entrega: 7 de diciembre de 2026}

\begin{document}
\maketitle

\begin{abstract}
\noindent Se propone predecir, \textbf{en el momento en que un perro o un gato ingresa a un refugio},
cuál será su desenlace entre cinco clases: adopción, transferencia a otra organización, devolución a
su dueño, eutanasia o muerte. Se usan los registros públicos de ingresos y egresos del
\emph{Austin Animal Center} (Texas), con más de 160\,000 estancias de perros y gatos entre 2013 y 2025.
El problema es multiclase, \textbf{fuertemente desbalanceado} (la clase más rara representa
$\approx$1\,\%), con variables categóricas de alta cardinalidad (más de mil razas y cientos de colores),
requiere unir dos tablas por animal y fecha, y se evalúa con una \textbf{partición temporal} (entrenar
con años anteriores y probar con los más recientes). Se compara un circuito variacional con
\emph{data re-uploading} contra modelos clásicos de capacidad similar. No encontramos trabajos de
aprendizaje automático cuántico sobre este problema.
\end{abstract}

\tableofcontents
\bigskip

%=====================================================================
\section{Motivación del cambio y justificación de la complejidad}
%=====================================================================
La propuesta anterior (perros \emph{vs.}\ gatos) era binaria y ya había sido resuelta con un enfoque
casi idéntico en \cite{martin2026}. Esta propuesta conserva el contexto de los refugios, pero cambia
a un problema \textbf{multiclase} sobre \textbf{datos tabulares heterogéneos}, distinto del caso de
imágenes. Tiene menos clases que la Opción 1 (5 frente a 7), pero es más difícil en casi todas las
demás dimensiones, como resume la tabla~\ref{tab:complejidad}.

\begin{table}[H]
\centering\small
\caption{Comparación de complejidad con la propuesta de clase (Opción 1).}
\label{tab:complejidad}
\begin{tabular}{@{}p{4.6cm}p{4.4cm}p{5.8cm}@{}}
\toprule
Aspecto & Opción 1: Dry Bean & Esta propuesta: refugio de Austin \\ \midrule
Número de clases & 7 & 5 \\
Desbalance (clase mayor / menor) & $\approx 6.8:1$ & $\approx \mathbf{50:1}$ \\
Tamaño & 13\,611 granos & $\approx$160\,000 estancias \\
Tipo de variables & 16 numéricas, ya limpias & mixtas: categóricas de alta cardinalidad
(razas, colores, tipo y condición de ingreso), numéricas (edad) y temporales \\
Preparación de datos & ninguna & unir ingresos y egresos por animal y fecha; animales con varias
estancias; limpieza de edades y razas \\
Riesgo de fuga & bajo & alto: variables que solo existen al egreso (estado reproductivo tras
esterilización, duración de la estancia) deben excluirse \\
Partición & aleatoria estratificada & \textbf{temporal}, con cambio de distribución entre años
(p.\ ej.\ pandemia 2020) \\
Codificación cuántica & valores continuos directos & primero hay que convertir categorías en
valores numéricos sin fuga \\
\bottomrule
\end{tabular}
\end{table}

%=====================================================================
\section{Narrativa y problema real}
%=====================================================================
Los refugios municipales operan con espacio, personal y presupuesto limitados. Cuando un animal
ingresa, el personal decide qué recursos dedicarle: promoción para adopción, búsqueda activa del
dueño, contacto con organizaciones de rescate o atención médica prioritaria. Saber desde el ingreso
qué animales tienen mayor riesgo de un desenlace negativo (eutanasia o muerte) permite intervenir a
tiempo, por ejemplo con hogares temporales, campañas dirigidas o atención veterinaria temprana.

El sistema propuesto es una \textbf{herramienta de apoyo}, no de decisión: su objetivo es priorizar
ayuda, nunca justificar una eutanasia. Por eso la métrica principal no es la \emph{accuracy} sino la
capacidad de detectar las clases minoritarias.

\subsection{Pregunta de investigación y objetivos}
\begin{quote}
\emph{¿Puede un clasificador variacional cuántico de pocos qubits predecir el desenlace de un animal
desde su ingreso, en un problema multiclase muy desbalanceado y con cambio temporal de distribución,
con un desempeño comparable al de modelos clásicos de capacidad similar?}
\end{quote}
Objetivos específicos:
\begin{enumerate}[label=(\roman*),leftmargin=*]
\item construir el conjunto de datos uniendo ingresos y egresos, usando solo información disponible
al ingreso;
\item diseñar una codificación de las variables categóricas en un vector numérico pequeño, sin fuga de
información, apta para codificarse en ángulos de rotación;
\item implementar y entrenar un clasificador con \emph{data re-uploading} y una pérdida que compense
el desbalance;
\item compararlo con modelos clásicos mediante F1 macro, \emph{balanced accuracy}, sensibilidad por
clase y matriz de confusión;
\item analizar profundidad, re-subida de datos, \emph{shots}, costo y estabilidad del desempeño año
con año;
\item discutir con honestidad la existencia o ausencia de ventaja cuántica.
\end{enumerate}

\subsection{Formalización}
Sea $\mathcal{D}=\{(\bm{r}_i,y_i,t_i)\}_{i=1}^{N}$, donde $\bm{r}_i$ es el registro de ingreso de la
estancia $i$, $t_i$ su fecha y $y_i\in\{0,\dots,4\}$ el desenlace. Una transformación
$\phi:\bm{r}\mapsto\bm{x}\in\R^{d}$ (ajustada solo con entrenamiento) produce las variables numéricas.
Se busca $f_{\bm\theta}:\R^{d}\to\Delta^{4}$ que minimice la entropía cruzada \textbf{ponderada por
clase}
\begin{equation}
\mathcal{L}(\bm\theta)=-\frac{1}{\sum_i w_{y_i}}\sum_{i=1}^{N} w_{y_i}\,\log\big[f_{\bm\theta}(\bm{x}_i)\big]_{y_i},
\qquad w_c=\frac{N}{5\,N_c},
\label{eq:loss}
\end{equation}
donde $N_c$ es el número de ejemplos de la clase $c$ en entrenamiento. Sin los pesos $w_c$, un modelo
que nunca predice eutanasia ni muerte tendría buena \emph{accuracy} y sería inútil. El problema es
\textbf{no convexo} y se reportan resultados con tres semillas.

%=====================================================================
\section{Datos}
%=====================================================================
Se usan dos conjuntos públicos (dominio público) del portal de datos abiertos de la Ciudad de Austin:
\emph{Austin Animal Center Intakes} (173\,812 ingresos) \cite{austinintakes} y \emph{Austin Animal
Center Outcomes} \cite{austinoutcomes}, ambos del 1 de octubre de 2013 al 5 de mayo de 2025. Se
restringe el estudio a perros y gatos.

\paragraph{Variables de ingreso.} Tipo de ingreso (callejero, entrega por el dueño, asistencia
pública, abandono, etc.), condición al ingreso (normal, herido, enfermo, lactante, neonato, edad
avanzada, etc.), especie, sexo y estado reproductivo al ingreso, edad al ingreso, raza, color, si
tiene nombre y fecha de ingreso.

\paragraph{Clases.} A partir de los egresos de perros y gatos (conteos aproximados del portal):
\begin{table}[H]
\centering\small
\begin{tabular}{@{}lrr@{}}
\toprule
Desenlace (clase) & Estancias & Proporción \\ \midrule
Adopción & 83\,256 & 50.8\,\% \\
Transferencia a otra organización & 47\,268 & 28.9\,\% \\
Devolución al dueño (incluye \emph{Rto-Adopt}) & 26\,794 & 16.4\,\% \\
Eutanasia & 4\,778 & 2.9\,\% \\
Muerte (incluye \emph{Disposal}) & 1\,669 & 1.0\,\% \\
\bottomrule
\end{tabular}
\end{table}
\noindent Se descartan los desenlaces sin significado para el problema (extraviado, robado,
reubicado; menos de 0.2\,\%) y los registros sin desenlace.

\paragraph{Construcción del conjunto.} Un mismo animal puede tener varias estancias. Cada egreso se
empareja con el ingreso más reciente del mismo \texttt{Animal ID} anterior a él; se descartan los
pares inconsistentes (egreso sin ingreso previo). Además, se agrega una variable con el número de
estancias previas del animal, calculada solo con el pasado.

\paragraph{Partición temporal.} Entrenamiento: ingresos de octubre de 2013 a diciembre de 2021;
validación: 2022; prueba: enero de 2023 a mayo de 2025. Así se evalúa la generalización hacia el
futuro, que es el uso real del sistema, y se mide la degradación entre años.

\paragraph{Variables excluidas por fuga.} El sexo al egreso (los animales se esterilizan antes de ser
adoptados), la fecha y la duración de la estancia, y el subtipo de desenlace, porque no se conocen al
ingreso.

%=====================================================================
\section{Preprocesamiento: de categorías a ángulos}
%=====================================================================
Las variables categóricas no pueden codificarse directamente en un circuito pequeño: \emph{one-hot}
sobre más de mil razas requeriría miles de dimensiones. Se proponen los siguientes pasos, todos
ajustados solo con entrenamiento:
\begin{enumerate}[leftmargin=*]
\item \textbf{Limpieza y agrupación.} Edad en días (escala logarítmica); raza dividida en
indicador de mezcla, raza principal y grupo de raza; color reducido a color base y patrón;
agrupación de categorías raras (menos de 0.5\,\%) en ``otra''.
\item \textbf{Codificación por objetivo con validación cruzada} (\emph{out-of-fold target encoding}):
cada categoría de alta cardinalidad se reemplaza por la frecuencia suavizada de cada desenlace en esa
categoría, calculada con particiones internas del entrenamiento para que ningún registro vea su propia
etiqueta.
\item \textbf{Variables cíclicas}: mes de ingreso como $(\sin,\cos)$.
\item \textbf{Estandarización y PCA} a $d\in\{6,8,12\}$ componentes, con reporte de la varianza retenida.
\item \textbf{Escalado a $[0,\pi]$} por min--max ajustado en entrenamiento.
\end{enumerate}

\paragraph{Submuestreo para el simulador.} Simular el circuito sobre $\approx110\,000$ estancias de
entrenamiento por época es costoso. Se entrenará con un subconjunto estratificado (p.\ ej.\ 20\,000
estancias, con todas las de las clases minoritarias), documentado y usado también para los baselines
clásicos de capacidad comparable. Validación y prueba se evalúan completos.

%=====================================================================
\section{Enfoque cuántico}
%=====================================================================
\subsection{Fundamento: \emph{data re-uploading}}
Si los datos se codifican una sola vez con compuertas $e^{-ixH/2}$, la salida del circuito es una
serie de Fourier con frecuencias $\{-1,0,1\}$ por variable \cite{schuld2021}. El \emph{data
re-uploading} \cite{perez2020} intercala $L$ capas de codificación y bloques entrenables,
\begin{equation}
U_{\bm\theta}(\bm{x})=\prod_{l=L}^{1}W_l(\bm\theta_l)\,S_l(\bm{x};\bm\alpha_l),\qquad
S_l(\bm{x};\bm\alpha_l)=\bigotimes_{q=1}^{n}R_X(\alpha_{l,q}\,x_q),
\end{equation}
con lo que cada capa amplía el espectro de frecuencias. $W_l$ aplica $\mathrm{Rot}(\phi,\vartheta,\omega)$
por qubit y un anillo de CNOT. Con $d=8$ se usan $n=8$ qubits, o $n=4$ qubits con dos variables por
qubit ($R_Y$ y $R_Z$), y se comparan ambas opciones.

\subsection{Lectura multiclase}
Se miden los primeros 3 qubits ($2^3=8\geq5$ estados) y se asigna un estado base $\ket{s_c}$ a cada
clase; las probabilidades se transforman con un \emph{softmax} de temperatura entrenable $\beta$:
\begin{equation}
\big[f_{\bm\theta}(\bm{x})\big]_c=\frac{\exp\!\big(\beta\,p_{s_c}(\bm{x})\big)}
{\sum_{c'}\exp\!\big(\beta\,p_{s_{c'}}(\bm{x})\big)}.
\end{equation}
Con $n=8$ y $L=4$ capas el modelo tiene $(3+1)\cdot8\cdot4+1=129$ parámetros.

\begin{figure}[H]
\centering
\begin{tikzpicture}[node distance=0.5cm and 0.8cm, >=Stealth,
  box/.style={draw, rounded corners, align=center, font=\small, minimum height=0.9cm, inner sep=4pt}]
\node[box] (in) {Ingresos\\+ egresos};
\node[box, right=of in] (join) {Unión por\\animal y fecha};
\node[box, right=of join] (enc) {Codificación\\por objetivo};
\node[box, right=of enc] (pca) {PCA\\$[0,\pi]$};
\node[box, right=of pca] (qc) {Circuito\\re-uploading};
\node[box, right=of qc] (out) {5\\desenlaces};
\draw[->] (in) -- (join);
\draw[->] (join) -- (enc);
\draw[->] (enc) -- (pca);
\draw[->] (pca) -- (qc);
\draw[->] (qc) -- (out);
\end{tikzpicture}
\caption{Pipeline propuesto. Todas las etapas que dependen de los datos se ajustan solo con el
periodo de entrenamiento.}
\end{figure}

\subsection{Optimización}
Adam (tasa inicial $0.05$, decaimiento coseno, \emph{minibatch} de 128, 20 épocas) sobre simulador de
vector de estado con diferenciación automática; selección del modelo por F1 macro en validación (2022).

\subsection{¿Por qué puede resolver el problema?}
\begin{enumerate}[leftmargin=*]
\item \textbf{No linealidad.} Los desenlaces dependen de interacciones (por ejemplo, un gato neonato
callejero en verano tiene un riesgo muy distinto al de un perro adulto entregado por su dueño); las
re-subidas y el entrelazamiento producen términos cruzados entre variables.
\item \textbf{Compacidad.} Las interacciones se expresan con $\approx130$ parámetros, del orden de un
MLP pequeño, lo que hace interesante la comparación justa.
\item \textbf{Desbalance.} La pérdida ponderada \eqref{eq:loss} y la temperatura $\beta$ evitan que el
circuito colapse a la clase mayoritaria.
\item \textbf{Universalidad con matices.} Con suficientes capas la familia aproxima funciones acotadas
\cite{perez2020}, pero eso no garantiza que el optimizador las encuentre; la evidencia empírica es
parte central de la justificación.
\end{enumerate}

\subsection{Verificación y costo en hardware}
Reconstrucción independiente del circuito en Qiskit para comprobar las probabilidades,
transpilación a compuertas nativas (profundidad y número de CNOT) y estimación del costo en
\emph{shots} de entrenar con \emph{parameter-shift}.

%=====================================================================
\section{Librerías y justificación de su uso}
%=====================================================================
\begin{itemize}[leftmargin=*]
\item \textbf{pandas}: unión de las tablas de ingresos y egresos, limpieza y variables temporales.
\item \textbf{scikit-learn}: \emph{target encoding} con validación cruzada (\texttt{TargetEncoder}),
estandarización, PCA, baselines (incluido \emph{gradient boosting}) y métricas.
\item \textbf{PennyLane}: definición del circuito y gradientes automáticos.
\item \textbf{Qiskit}: verificación cruzada y estimación de recursos.
\item \textbf{NumPy, matplotlib}: cálculo numérico y gráficas.
\end{itemize}
No se necesitan redes preentrenadas ni GPU.

%=====================================================================
\section{Baselines clásicos y diseño experimental}
%=====================================================================
\paragraph{Métricas.} F1 macro (principal), \emph{balanced accuracy}, sensibilidad (\emph{recall}) de
eutanasia y muerte, pérdida logarítmica, matriz de confusión y \emph{accuracy} como referencia. Media
y desviación estándar sobre tres semillas.

\paragraph{Baselines} (mismas particiones, mismas $d$ variables y mismos pesos de clase):
\begin{itemize}[leftmargin=*]
\item clasificador que siempre predice la clase mayoritaria y clasificador de centroides;
\item regresión logística multinomial;
\item MLP con $\approx130$ parámetros y MLP más grande;
\item Random Forest y \emph{gradient boosting} (\texttt{HistGradientBoostingClassifier});
\item \emph{gradient boosting} con todas las variables sin PCA, como cota superior.
\end{itemize}

\paragraph{Ablaciones.}
(a) profundidad $L\in\{1,2,4,6\}$;
(b) con y sin re-subida de datos;
(c) $d\in\{6,8,12\}$ y una o dos variables por qubit;
(d) con y sin pesos de clase;
(e) desempeño por año de prueba (2023, 2024, 2025), para medir el efecto del cambio temporal;
(f) muestreo finito con $S\in\{100,500,2\,000,10\,000\}$ \emph{shots}.

%=====================================================================
\section{Resultados esperados y criterio de éxito}
%=====================================================================
Se espera que \emph{gradient boosting} sobre todas las variables sea el mejor modelo, y que la
eutanasia y la muerte sean las clases más difíciles. Se considera exitoso el proyecto si:
\begin{itemize}[leftmargin=*]
\item el clasificador cuántico supera claramente a la clase mayoritaria y a los centroides en F1 macro
y \emph{balanced accuracy};
\item queda dentro de unos pocos puntos de F1 macro de un MLP de capacidad similar;
\item detecta una fracción útil de los casos de eutanasia y muerte;
\item se cuantifica la degradación entre años y con pocos \emph{shots}.
\end{itemize}
Un resultado negativo también es válido y se discutirá como tal.

%=====================================================================
\section{Riesgos, ética y mitigación}
%=====================================================================
\begin{itemize}[leftmargin=*]
\item \textbf{Fuga de información}: lista explícita de variables excluidas y verificación de que cada
variable se conoce al ingreso.
\item \textbf{Clases muy raras}: pérdida ponderada, métricas macro y análisis por clase.
\item \textbf{Cambio de distribución}: partición temporal y reporte por año.
\item \textbf{Uso indebido}: el modelo se plantea solo para priorizar ayuda; se discutirán sesgos por
raza (p.\ ej.\ razas estigmatizadas) en el análisis de errores.
\item \textbf{Tiempo de simulación}: submuestreo estratificado documentado y circuitos de 4--8 qubits.
\end{itemize}

%=====================================================================
\section{Cronograma}
%=====================================================================
\begin{table}[H]
\centering\small
\begin{tabular}{@{}p{3.0cm}p{9.2cm}p{2.3cm}@{}}
\toprule
Fechas & Actividad & Hito \\ \midrule
17--18 nov & Validación con la profesora; investigación teórica; formalización. & Validación \\
19--22 nov & Descarga, unión de tablas, limpieza y EDA; codificación de categorías; baselines
clásicos; primer circuito; diapositivas. & Avances de código \\
23 nov & Presentación de pre-entrega. & \textbf{Pre-entrega} \\
24--29 nov & Circuito en PennyLane, verificación con Qiskit, entrenamiento con tres semillas. & Modelo entrenado \\
30 nov--3 dic & Ablaciones, análisis por año, muestreo finito y costo. & Resultados \\
4--6 dic & Cuaderno final reproducible, conclusiones y revisión contra la rúbrica. & Cierre \\
7 dic & Entrega del cuaderno y del documento. & \textbf{Entrega} \\
\bottomrule
\end{tabular}
\end{table}

%=====================================================================
\begin{thebibliography}{9}
\bibitem{austinintakes} City of Austin. \emph{Austin Animal Center Intakes (10/01/2013 to
05/05/2025)}. Portal de datos abiertos, dominio público.
\url{https://data.austintexas.gov/d/wter-evkm}.
\bibitem{austinoutcomes} City of Austin. \emph{Austin Animal Center Outcomes (10/01/2013 to
05/05/2025)}. Portal de datos abiertos, dominio público.
\url{https://data.austintexas.gov/d/9t4d-g238}.
\bibitem{martin2026} D.\ Martín-Pérez, F.\ Rodríguez-Díaz, D.\ Gutiérrez-Avilés, A.\ Troncoso y
F.\ Martínez-Álvarez. \emph{Hybrid Classical-Quantum Transfer Learning with Noisy Quantum Circuits}.
Computer Modeling in Engineering \& Sciences, 148(1), 2026.
\bibitem{perez2020} A.\ Pérez-Salinas, A.\ Cervera-Lierta, E.\ Gil-Fuster y J.\ I.\ Latorre.
\emph{Data re-uploading for a universal quantum classifier}. Quantum, 4:226, 2020.
\bibitem{schuld2021} M.\ Schuld, R.\ Sweke y J.\ J.\ Meyer. \emph{Effect of data encoding on the
expressive power of variational quantum-machine-learning models}. Phys.\ Rev.\ A, 103:032430, 2021.
\bibitem{bergholm2018} V.\ Bergholm \emph{et al.} \emph{PennyLane: Automatic differentiation of hybrid
quantum-classical computations}. arXiv:1811.04968, 2018.
\end{thebibliography}

\end{document}
